\section{偏序集}

\begin{dfn} \textbf{偏序关系, 偏序集} Partial Order, Poset\\
若集合$X$上的关系$R$满足自反性、反对称性、传递性，称$R$为偏序关系, $(X,R)$称为偏序集. 
\end{dfn}
\par 当$R$定义明确, 或不强调关系的具体定义时, 往往略去不写. 以下用$\le$代表特定的偏序关系, 并记$x<y\iff x\le y\land x\neq y$. 对于偏序集$X$的子集$S$, 如无特殊说明, 则$S$中元素的偏序关系与$X$一致. 

\begin{dfn} \textbf{区间} Interval\\
设$X$是偏序集, $a,b\in X$, $a\le b$, 则$\{x\in X\vert a\le x\le b\}$称为闭区间, 记作$[a,b]$;  $\{x\in X\vert a<x<b\}$称为开区间, 记作$(a,b)$. 
\end{dfn}
类似可定义半开半闭区间$[a,b)$或$(a,b]$. 

\begin{dfn} \textbf{局部有限偏序集} Locally Finite Poset\\
设$X$是偏序集, 若对任意$a,b\in X$, 当$a\le b$时, 闭区间$[a,b]$是有限集, 称$X$为局部有限偏序集. 
\end{dfn}

\begin{dfn} \textbf{全序关系, 全序集} Total Order\\
若集合$X$上的偏序关系$R$满足$\forall x,y\in X: xRy \lor yRx$, 称$R$为全序关系, $(X,R)$称为全序集. 
\end{dfn}
上述定义中的性质可表达为$R\cup R^{-1}=X\times X$.

\begin{dfn} \textbf{链, 反链} Chain, Antichain\\
设$X$是偏序集, $A\subseteq X$. 若$A$是全序集, 称$A$为$X$的一个链. 若$\forall x,y\in A, x\neq y: \lnot (x\le y \lor y\le x)$, 称$A$为$X$的一个反链. 
\end{dfn}


\begin{dfn} \textbf{最小元，最大元} Smallest/Largest Element\\
设$X$是偏序集, 对$a\in X$, 若$\forall x\in X: a\le x$, 称$a$是最小元; 若$\forall x\in X: x\le a$, 称$a$是最大元. 
\end{dfn}
根据反对称性, 偏序集的最小（大）元若存在则唯一. 

\begin{dfn} \textbf{极小元，极大元} Minimal/Maximal Element\\
设$X$是偏序集, 对$a\in X$, 若$x\le a\Rightarrow x=a$, 称$a$是极小元; 若$a\le x\Rightarrow x=a$, 称$a$是极大元. 
\end{dfn}

\begin{dfn} \textbf{下界，上界} Lower/Upper Bound\\
设$X$是偏序集, $A\subseteq X$, $b\in X$, 若$\forall x\in A: b\le x$, 称$b$是$A$的一个下界; 若$\forall x\in A: x\le b$, 称$b$是$A$的一个上界. 
\end{dfn}
若$b\in A$是$A$的一个下界, 则$b$是$A$的最小元.

\begin{dfn} \textbf{下确界，上确界} Infimum, Supremum\\
设$X$是偏序集, $A\subseteq X$, 定义$A$的下界的集合$A_*=\{b\in X\vert \forall x\in A: b\le x\}$. 若$A_*$有最大元, 称之为$A$的下确界, 记作$\text{inf} A$. 定义$A$的上界的集合$A^*=\{b\in X\vert \forall x\in A: x\le b\}$. 若$A_*$有最小元, 称之为$A$的上确界, 记作$\text{sup} A$. 
\end{dfn}
\par 下确界、上确界若存在, 则唯一. 

\begin{dfn} \textbf{偏序集的覆盖关系} Covering Relation\\
设$X$是偏序集, $X$上的关系$<_c$定义为$x<_c y \iff x<y \land (\lnot \exists z: x<z<y)$. $<_c$称为偏序集$X$的覆盖关系. 
\end{dfn}

\par 基于覆盖关系可绘制偏序集$X$的Hasse图: 图中每个节点代表$X$的一个元素, 当$x<_c y$时$x,y$连边, 并将$x$置于$y$的下方. 

\begin{dfn} \textbf{偏序集的相似} Similarity\\
设$X,Y$是偏序集, 若存在双射$f:X\to Y$, 满足$\forall a,b\in X$, 在$X$中$a\le b$当且仅当在$Y$中$f(a)\le f(b)$, 称$X$与$Y$相似, 记作$X\cong Y$. 
\end{dfn}
上述保序条件的一个等价表述为: $\forall a,b\in X$, 在$X$中$a<b$当且仅当在$Y$中$f(a)<f(b)$. 这一条件事实上蕴含$f$是单射. 给定全集$E$, 相似关系是$\mathcal{P}(E)$中的偏序集之间的等价关系. 

\begin{thm} \textbf{Zorn引理} Zorn's Lemma\\
设$X$是偏序集, 若对任意链$A\subseteq X$, $X$中存在$A$的上界, 则$X$有极大元. 
\end{thm}
\begin{proof}
\par 由条件知$X$非空. 设$\mathcal{X}=\{A\subseteq X| A\text{是链}\}$, 则由$\varnothing\in \mathcal{X}$知$\mathcal{X}$非空. 考虑偏序集$(\mathcal{X}, \subseteq)$, 若$\mathcal{X}$存在极大元$A$, 则存在$A$的上界$x_a\in X$, 满足$\forall a\in A: a\le x_a$. 设$x_a\le x$, 若$x_a < x$, 则$x\notin A$, 因此$A\subset A\cup \{x\}$, 且$A\cup \{x\}$是链, 与$A$是$\mathcal{X}$的极大元矛盾. 因此有$x_a=x$, 这推出$x_a$是$X$的极大元. 因此只需证$\mathcal{X}$中存在极大元. 
\par 由选择公理, 存在函数$f:\mathcal{P}(X)-\{\varnothing\}\to X$, 对$X$的任意非空子集$A$有$f(A)\in A$. 对任意$A\in \mathcal{X}$, 令$\hat{A}=\{x\in X\vert A\cup \{x\} \in \mathcal{X}\}$. 定义函数$g:\mathcal{X}\to \mathcal{X}$: 若$\hat{A}-A=\varnothing$, 令$g(A)=A$; 否则令$g(A)=A\cup f(\hat{A}-A)$. 注意到若$\hat{A}-A=\varnothing$, 则$A$是$\mathcal{X}$的极大元. 事实上, 若存在$A'\in \mathcal{X}$使得$A\subset A'$, 则存在$a\in A'-A$, 此时$A\cup \{a\}$也是链, 故$a\in \hat{A}-A$, 与$\hat{A}-A=\varnothing$矛盾. 因此只需证存在$A\in \mathcal{X}$使得$g(A)=A$. 
\par 若$\mathcal{X}$的子集$\mathcal{T}$满足:
\begin{enumerate}
    \item $\varnothing\in \mathcal{T}$;
    \item $\forall A\in \mathcal{T}: g(A)\in \mathcal{T}$;
    \item 若$\mathcal{C} \subseteq \mathcal{T}$是链, 则$\bigcup_{A\in \mathcal{C}} A \in \mathcal{T}$.
\end{enumerate}
称$\mathcal{T}$是塔. 易知$\mathcal{X}$是塔. 设$\mathcal{T}_0$是所有塔构成的集合的交, 则$\mathcal{T}_0$也是塔. 往证$\mathcal{T}_0$是链. 对$C\in \mathcal{T}_0$, 若$\forall A\in \mathcal{T}_0: A\subseteq C \lor C\subseteq A$, 则称$C$是可比的. 显然$\varnothing$是可比的. 设$C$是可比的, 记$\mathcal{U}_C=\{A\in \mathcal{T}_0\vert A\subseteq C \lor g(C)\subseteq A\}$, 则$\mathcal{U}_C$是塔. 事实上, (1) $\varnothing\in \mathcal{U}_C$; (2)若$A\in \mathcal{U}_C$, 分三种情况讨论: (i)$A\subset C$, 则$g(A)\subseteq C$或$C\subset g(A)$. 由于$g(A)-A$至多包含一个元素, 后一种情形是不可能的, 故$g(A)\subseteq C$. (ii)$A=C$, 则$g(C)=g(A)$. (iii) $g(C)\subseteq A$, 则由$A\subseteq g(A)$, $g(C)\subseteq g(A)$. 因此三种情况均有$g(A)\in \mathcal{U}_C$. (3) 若$\mathcal{C} \subseteq \mathcal{U}_C$是链, 首先$\bigcup_{A\in \mathcal{C}} A \in \mathcal{T}_0$; 若存在$A\in \mathcal{C}$, $g(C)\subseteq A$, 则$g(C)\subseteq \bigcup_{A\in \mathcal{C}} A$; 否则$\forall A\in \mathcal{C}$, $A\subseteq C$, 则$\bigcup_{A\in \mathcal{C}} A\subseteq C$. 故$\bigcup_{A\in \mathcal{C}} A \in \mathcal{U}_C$. 这证明了$\mathcal{U}_C$是塔. 由$\mathcal{T}_0$的定义, $\mathcal{T}_0\subseteq \mathcal{U}_C$; 又由于$\mathcal{U}_C\subseteq\mathcal{T}_0$, 得$\mathcal{T}_0=\mathcal{U}_C$, 即$\forall A\in \mathcal{T}_0$, $A\subseteq C \lor g(C)\subseteq A$. 而$A\subseteq C$蕴含$A\subseteq g(C)$, 这推出$g(C)$是可比的. 又由于$\mathcal{T}_0$中可比集合构成的链的并也是可比的, 知$\mathcal{T}_0$中可比集合的集合是塔, 从而是$\mathcal{T}_0$本身. 因此$\mathcal{T}_0$是链.
\par 取$\bar{A}=\bigcup_{A\in \mathcal{T}_0} A$, 则$\bar{A}\in \mathcal{T}_0$, 进而$g(\bar{A})\in \mathcal{T}_0$, 因此有$g(\bar{A})\subseteq \bar{A}$. 又由于$\bar{A}\subseteq g(\bar{A})$, 故$g(\bar{A})=\bar{A}$. 证毕. 
\end{proof}



\section{良序集}

\begin{dfn} \textbf{良序关系} Well Order\\
设$(X,R)$是偏序集, 若$X$的任意非空子集有最小元, 称$R$为良序关系, $(X,R)$称为良序集. 
\end{dfn}
\par 设$X$是良序集, 对任意$x_1,x_2\in X$, 令$A=\{x_1,x_2\}$, 知$x_1\le x_2 \lor x_2\le x_1$. 因此良序关系是全序关系. 

\begin{dfn} \textbf{良序集的延续} Continuation\\
设$A$是良序集, 若$B\subseteq A$, $\forall b\in B\forall a\in A-B: b<a$, 称$A$是良序集$B$的延续. 
\end{dfn}
\par 延续是良序集之间的偏序关系. 事实上, 自反性和反对称性显然. 若良序集$A$是$B$的延续, $B$是$C$的延续, 设$c\in C, a\in A-C$, 若$a\in B$, 则$a\in B-C$; 否则$a\in A-B$, $c\in B$; 两种情况均有$c<a$. 这表明$A$是$C$的延续.

\begin{thm}\textbf{超限归纳法原理} Principle of Transfinite Induction\\
设$X$是良序集, $S\subseteq X$. 若对任意$x\in X$, $(\forall s<x: s\in S) \Rightarrow x\in S)$, 则$S=X$. 
\label{thm:transinduct}
\end{thm}
\begin{proof}
若$X-S$非空, 由$X$是良序集, $X-S$存在最小元$x_l$, 满足$x<x_l \Rightarrow x\in S$, 由条件知$x_l \in S$, 矛盾. 因此$X-S=\varnothing$, $X=S$.
\end{proof}

\begin{thm} \textbf{良序定理} Well Ordering Theorem\\
设$X$是集合, 存在$X$上的良序关系. 
\end{thm}
\begin{proof}
令$\mathcal{W}=\{(A,\le)\in \mathcal{P}(X)\times \mathcal{P}(X\times X)\vert \le\text{是}A\text{上的良序关系}\}$, 由$\varnothing\in \mathcal{W}$知$\mathcal{W}$非空. $\mathcal{W}$在延续关系下构成偏序集. 设$\mathcal{C}\subseteq \mathcal{W}$是链, 则存在$U=\bigcup_{C\in \mathcal{C}} C$上的良序关系, 使$\forall C\in \mathcal{C}$, $U$是$C$的延续. 事实上, 对$c_1,c_2\in U$, 存在$C_1,C_2\in \mathcal{C}$, $c_1\in C_1$, $c_2\in C_2$; 又由于$\mathcal{C}$是链, $C_1,C_2$存在延续关系. 不妨设$C_1$是$C_2$的延续, 则$c_2\in C_1$. 定义$c_1,c_2$在$U$中的关系与$C_1$一致, 易知$\forall C\in \mathcal{C}$, 若$c_1,c_2\in C$, 则$c_1,c_2$在$C$中的关系与$C_1$中一致, 因此$U$上的上述关系是良定义的, 且易知这一关系仍是偏序关系. 设非空集合$A\subseteq U$, 则存在$C\in \mathcal{C}$, $A\cap C\neq \varnothing$. 设$m$是$A\cap C$的最小元, 则$m$也是$A$的最小元, 因此$U$是良序集. 对$C\in \mathcal{C}$, 设$c\in C$, $u\in U-C$, 则存在$D\in \mathcal{C}$, $D\neq C$, $u\in D$, 知$D$是$C$的延续, 且$c<u$. 故$U$是$C$的延续.
\par 由Zorn引理, $\mathcal{W}$存在极大元$M$. 我们证明$M=X$. 事实上, 若$x\in X-M$, 令$M'=M\cup \{x\}$, 并补充定义$m\le x, \forall m\in M'$, 则$M'$仍是良序集（对非空集合$A\subset M'$, 若$A\cap M$非空, $A\cap M$在$M$中的最小元即为$A$的最小元; 否则$A=\{x\}$, $x$是最小元）, 且是$M$的延续, 与$M$为极大元矛盾. 
\end{proof}

\begin{thm} \textbf{超限递归定理} Transfinite Recursion Theorem\\
设$W$是良序集, 对$a\in W$, 令$s(a)=\{x\in W\vert x<a\}$. 给定函数$f:\{t\in X^{s(a)}\vert a\in W\}\to X$, 存在唯一的函数$U:W\to X$, 满足$\forall a\in W: U(a)=f(U|_{s(a)})$.
\end{thm}
\begin{proof}
唯一性利用超限归纳法即证, 只需证存在性. 称$A\subseteq W\times X$是$f$封闭的, 如果$\forall a\in W \forall t\in X^{s(a)}: (\forall c\in s(a): (c,t(c))\in A )\Rightarrow (a,f(t))\in A$. $W\times X$是$f$封闭的. 令$U$为$W\times X$的$f$封闭的子集交集, 知$U$是$f$封闭的. 以下用超限归纳法证明$U$是函数, 记$S=\{a\in W\vert (a,x)\in U\land (a,y)\in U\Rightarrow x=y\}$, 往证$s(a)\subseteq  S\Rightarrow a\in S$. 
\par 当$s(a)\subseteq S$, $t=\{(c,x)\in U\vert c<a\}$是$s(a)$上的函数. 由$U$是$f$封闭的, $(a,f(t))\in U$. 若存在$y\neq f(t)$, $(a,y)\in U$, 则$U-\{(a,y)\}$是$f$封闭的. 事实上, 设$b\in W$, $r\in X^{s(b)}$, $r\subseteq U-\{(a,y)\}$. 若$b=a$, 必有$r=t$, 此时$(a,f(t))\in U-\{(a,y)\}$; 若$b\neq a$, 亦有$(b,f(r))\in U-\{(a,y)\}$. 这与$U$的定义矛盾. 证毕. 
\end{proof}
\par 超限递归定理的应用称为超限递归定义(definition by transfinite induction).  

\begin{lem}\textbf{良序集到自身相似关系的性质}\\
设$X$是良序集, 存在函数$f:X\to X$满足$\forall a,b\in X$, $a\le b$当且仅当$f(a)\le f(b)$, 则有$\forall a\in X: a\le f(a)$. 
\label{lem:wossim}
\end{lem}
\begin{proof}
用反证法, 设存在最小的$b\in X$, $b>f(b)$. 一方面$\forall a<b: a\le f(a)$, 故$f(b)\le f(f(b))$; 另一方面由$f$的保序性$f(b)>f(f(b))$, 矛盾. 
\end{proof}

\begin{pro}\textbf{良序集的比较定理} Comparability Theorem for Well Ordered Sets\\
设$X,Y$是良序集, 则以下三种情况有且仅有一种发生: 
\begin{enumerate}
    \item $X \cong Y$;
    \item 存在$a\in X$, $Y\cong s(a)=\{x\in X\vert x<a\}$;
    \item 存在$b\in Y$, $X\cong s(b)=\{y\in Y\vert y<b\}$.
\end{enumerate}
\label{thm:woscomp}
\end{pro}
\begin{proof}
先证三种情况必有一种发生. 设情形(2)(3)均不成立, 往证$X\cong Y$. 定义函数$f:\{t\in Y^{s(a)}\vert a\in X\}\to Y$如下: 若存在$t(s(a))$的上界$b\in Y-t(s(a))$, 则取满足条件的最小元; 否则取$Y$的最小元. 由超限递归定理, 存在唯一的函数$U: X\to Y$, 满足$\forall a\in X: U(a)=f(U|_{s(a)})$. 采用超限归纳法证明$\forall a\in X: U|_{s(a)}$是$s(a)$到$s(U(a))$的双射. 设$\forall x<a: U|_{s(x)}$是$s(x)$到$s(U(x))$的双射, 我们有: \par (i)$\forall x<a: U(x)<U(a)$. 断言存在$U(s(a))$的上界$b\in Y-U(s(a))$, 从而$U(a)$是$U(s(a))$的上界. 若不然, 设存在$y\in Y-U(s(a))$, 则存在$x<a$, $U(x)>y$. 而$U|_{s(x)}$是$s(x)$到$s(U(x))$的双射, 故$y\in U(s(x))$, 矛盾. 因此$U(s(a))=Y$. 任取$x_1<x_2\in s(a)$, 由归纳假设有$U(x_1)\in s(U(x_2))$, 亦即$U(x_1)<U(x_2)$; 反之亦然. 因此$U: s(a)\to Y$是保序的双射, $s(a)\cong Y$, 与假设矛盾. 
\par (ii)$U|_{s(a)}: s(a)\to s(U(a))$是单射. 任取$x_1\neq x_2\in s(a)$, 不妨设$x_1<x_2$, 由归纳假设有$U(x_1)\in s(U(x_2))$, 亦即$U(x_1)<U(x_2)$. 
\par (iii)$U|_{s(a)}: s(a)\to s(U(a))$是满射. 若不然, 存在最小的$y\in s(U(a))$, $y\notin U(s(a))$. 若存在$x\in s(a)$, $U(x)>y$, 则与$U|_{s(x)}$为$s(x)$到$s(U(x))$双射矛盾. 故$y$是$U(s(a))$的一个上界, 这与$U(a)$的选取矛盾. 
\par 综上我们证明了$\forall a\in X: U|_{s(a)}$是$s(a)$到$s(U(a))$的双射, 由此容易验证$U:X\to Y$的保序性, 只需证明$U$是满射. 若不然, 取$Y-U(X)$的最小元$b$, 易知$b$是$U(X)$的上界, 从而$U(X)=s(b)$, $X\cong s(b)$, 矛盾. 因此我们证明了$U(X)=Y$, $X\cong Y$. 
\par 若(1)(2)同时发生, 则$X\cong s(a)$. 设$g: X\to s(a)$是保序的双射, 则$g(a)<a$, 与引理\ref{lem:wossim}矛盾. 同理(1)(3)不可能同时发生. 若(2)(3)同时发生, 记$f:X\to s(b)$, $g: Y\to s(a)$为保序的双射, 则$f\circ g: Y \to s(b)$是$Y$到自身的保序函数, $f(g(b))<b$同样与引理\ref{lem:wossim}矛盾.
\end{proof}

\begin{col}\textbf{良序集及其子集的相似关系}\\
设$X$是良序集, $A\subseteq X$, 则$A\cong X$, 或存在$a\in X$, $A\cong s(a)=\{x\in X\vert x<a\}$.
\end{col}
\begin{proof}
若不然, 由定理\ref{thm:woscomp}, 存在$a\in A$, $X\cong s(a)\cap A$, 这与引理\ref{lem:wossim}矛盾. 
\end{proof}

\section{序数}

\begin{dfn} \textbf{序数} Ordinal Number\\
设$\alpha$是良序集, 满足$\forall a\in \alpha: a=\{x\in \alpha: x<a\}$, 则称$\alpha$为序数.
\end{dfn}
自然数和自然数集都是序数. 作为序数的自然数集也记为$w$.  

\begin{col} \textbf{序数是传递集}\\
设$\alpha$是序数, 则$a\in \alpha \Rightarrow a\subseteq \alpha$. 
\end{col}

\begin{pro} \textbf{序数的元素是序数}\\
设$\alpha$是序数, $a\in \alpha$, 则$a$是序数. 
\end{pro}
\begin{proof}
由$a\subseteq \alpha$知$a$是良序集. 设$x\in a$, 则$x=\{s\in \alpha\vert s<x\}=\{s\in a\vert s<x\}$. 
\end{proof}

\begin{pro} \textbf{相似的序数相等}\\
设$\alpha, \beta$是序数, 则$\alpha \cong  \beta\Rightarrow \alpha =\beta$. 
\label{pro:osimeq}
\end{pro}
\begin{proof}
设$f:\alpha\to \beta$是保持序关系的双射, 用超限归纳法证明$\forall a\in \alpha: f(a)=a$. 记$s(a)=\{x\in \alpha: x<a\}$, 由保序性知$f(s(a))=\{y\in \beta: y<f(a)\}$. 设命题对$s(a)$中任意元素成立, 则$s(a)=f(s(a))$; 又由$\alpha, \beta$是序数, $a=s(a)$, $f(a)=\{y\in \beta: y<f(a)\}$, 因此$a=f(a)$. 这表明$f$是恒等映射, $\alpha=\beta$.
\end{proof}

\begin{dfn} \textbf{序数的序}\\
设$E$是序数的非空集合, 定义$E$上的小于关系$<$如下: 对$\alpha, \beta\in E$, $\alpha<\beta \iff \alpha\in \beta$.
\end{dfn}

\begin{thm} \textbf{序数的集合是良序集}\\
序数的非空集合是良序集.
\end{thm}
\begin{proof}
\par 首先验证序数间的$<$关系导出的$\le$是偏序关系. 若$\alpha<\beta$与$\beta<\alpha$同时成立, 则由序数是传递集, $\alpha\in \alpha$, 此时由序数的定义$\alpha=\{x\in \alpha\vert x<\alpha\}$, 知$\alpha \notin \alpha$, 矛盾. 因此$<$满足禁对称性, 从而$\le$满足反对称性. $<$的传递性(也即$\le$的传递性)可由序数是传递集导出. 
\par 其次证明对任意序数$\alpha, \beta$, $\alpha<\beta$, $\beta<\alpha$, $\alpha=\beta$中至少有一个成立, 从而序数的非空集合是全序集. 若$\alpha\cong \beta$, 由命题\ref{pro:osimeq}, $\alpha=\beta$. 否则由良序集的比较定理, (i)存在$a\in \alpha$, $\beta \cong s(a)=\{x\in \alpha\vert x<a\}$; 或(ii)存在$b\in \beta$, $\alpha \cong s(b)=\{y\in \beta \vert y<b\}$. 在(i)下, 由$a=s(a)$知$\beta \cong a$; 又由$a$是序数及命题\ref{pro:osimeq}, $\beta=a$, 从而$\beta\in \alpha$. 在(ii)下, 同理有$\alpha\in \beta$.
\par 最后证明任意非空的序数集合$E$有最小元, 从而任意非空的序数集合是良序集. 任取$\alpha \in E$, 不妨设$\alpha$不是$E$的最小元, 此时存在$\beta \in E$, $\beta\in \alpha$, 因此$\alpha \cap E$非空. 由于$\alpha$是良序集, $\alpha \cap E$存在最小元$\alpha_0$. 对任意$\beta \in E$, 若$\beta <\alpha$, 则$\beta \in \alpha \cap E$, $\alpha_0\le \beta$; 否则$\alpha_0<\alpha\le \beta$. 因此$\alpha_0$也是$E$的最小元. 
\end{proof}
由上述证明过程知序数的$<$关系满足三歧性, 即$\alpha<\beta$, $\beta<\alpha$, $\alpha=\beta$有且仅有一个成立. 事实上, 若$\alpha<\beta$, $\alpha=\beta$同时成立, 同样有$\alpha\in \alpha$, 导出矛盾. 

\begin{col} \textbf{序数序关系的等价条件}\\
设$\alpha, \beta$是序数, 则以下条件等价:
\begin{enumerate}
    \item $\alpha\in \beta$;
    \item $\alpha\subset \beta$;
    \item $\alpha\neq \beta$, 且$\beta$是$\alpha$的延续.
\end{enumerate}
\end{col}
\begin{proof}(1)$\to$(2): 略.  
\par (2)$\to$(3): 由良序集序关系的三歧性, $\alpha \in \beta$ 或 $\beta \in \alpha$.  若$\beta \in \alpha$, 则$\beta\subseteq \alpha$, 与条件矛盾. 因此$\alpha \in \beta$, 此时$\alpha=\{x\in \beta\vert x<\alpha\}$, 从而$\beta$是$\alpha$的延续. 
\par (3)$\to$(1): 取$\beta-\alpha$的最小元$b$, 则由良序集延续的定义$\alpha=\{x\in \beta\vert x<b\}$. 又由$\beta$是序数, $b=\{x\in \beta\vert x<b\}$. 因此$\alpha=b$. 
\end{proof}

\begin{col} \textbf{序数与自然数集的关系}\\
设$\alpha$为序数. 若$\alpha$为有限集, 则$\alpha \in w$; 若$\alpha$为无限集, 则$w\subseteq \alpha$.
\label{col:inford_nat}
\end{col}
\begin{proof}
由序数序关系的三歧性, $\alpha\in w$或$w\subseteq \alpha$. 对于前者, $\alpha$显然是有限集; 对于后者, 由命题\ref{pro:mininfset}知$\alpha$为无限集. 
\end{proof}

\begin{dfn}\textbf{后继序数, 极限序数} Successor/Limit Ordinal\\
设$\alpha$是非零序数, 若存在序数$\beta$, 使得$\alpha=\beta^+$, 则称$\alpha$是后继序数, 否则称为极限序数. 
\end{dfn}
无限序数也称为超限序数(transfinite ordinal), $w$是最小的超限序数, 也是最小的极限序数.

\begin{pro} \textbf{序数的并集是序数}\\
设$\mathcal{C}$是序数的非空集合, 则$\alpha=\bigcup_{C\in \mathcal{C}}C$是序数, 且是$\mathcal{C}$的上确界. 
\end{pro}
\begin{proof}
\par 由$\mathcal{C}$在良序集的延续关系下构成链, 由良序定理的证明过程知$\alpha$是良序集, 且$\forall C\in \mathcal{C}$, $\alpha$是$C$的延续. 对$a\in \alpha$, 存在$A\in \mathcal{C}$, $a\in A$. 由$A$是序数, $a=\{x\in A\vert x<a\}=\{x\in \alpha \vert x<a\}$. 故$\alpha$是序数. 
\par $\forall C\in \mathcal{C}$, 由$\alpha$是$C$的延续知$C\le \alpha$. 设$\beta$是$\mathcal{C}$的任一上界, 则$\forall C\in \mathcal{C}$, $C\subseteq \beta$, 故$\alpha \subseteq \beta$, $\alpha\le \beta$, 因此$\alpha$是$\mathcal{C}$的上确界. 
\end{proof}

\begin{pro} \textbf{序数的后继是序数}\\
设$\alpha$是序数, 在$\alpha^+=\alpha \cup \{\alpha\}$中补充定义$\forall x\in \alpha: x\le \alpha$, 则$\alpha^+$是序数. 
\end{pro}

\begin{col} \textbf{不存在所有序数的集合}\\
不存在所有序数的集合. 
\end{col}
\begin{proof}
设$\mathcal{C}$是所有序数的集合, 则$\alpha =\bigcup_{C\in \mathcal{C}}C$是序数, 且是所有序数的上界. 但$\alpha^+$是序数且$\alpha<\alpha^+$, 这导致矛盾. 
\end{proof}
\par 假定所有序数的集合存在, 上述矛盾称为Burali-Forti悖论.

\begin{thm} \textbf{计数定理} Counting theorem\\
设$X$是良序集, 则存在唯一的序数$\alpha$使得$X\cong \alpha$, 记作$\alpha=\operatorname{ord}X$. 
\end{thm}
\begin{proof}
\par 由相似的序数相等, 唯一性显然, 下证存在性. 先证一个引理: 对良序集$A$, 若对任意$a\in A$, 存在序数$\alpha \cong s(a)=\{x\in A\vert x<a\}$, 则存在序数$\bar{\alpha}\cong A$. 事实上, 令$S(a)=\{\alpha\vert \alpha\text{是序数} \land \alpha\cong s(a)\}$, 则$S(a)$存在(是单件). 由替换公理模式, 存在$A$上的函数$F$, 满足$\forall a\in A: F(a)\text{是序数} \land F(a)\cong s(a)$. 令$\bar{\alpha}$为$F$的值域, 则$\bar{\alpha}$为良序集. 我们有:
\par (i) $F:A\to \bar{\alpha}$是保序的双射, 因此$A\cong \bar{\alpha}$. $F$显然是满射. 设$a,b\in A$, 若$a<b$, 则$s(a)\subset s(b)$. 设$f:s(a)\to F(a)$, $g:s(b)\to F(b)$是保序双射, 则$g\vert_{s(a)}$是$s(a)$到$\{x\in F(b)\vert x<g(a)\}=g(a)$的保序双射. 因此$g\vert_{s(a)}\circ f^{-1}$是$F(a)$到$g(a)$的保序双射, 从而$F(a)=g(a)\in F(b)$, 亦即$F(a)<F(b)$. 同理可证$a>b\Rightarrow F(a)>F(b)$. 
\par (ii) $\bar{\alpha}$是序数. 对任意$b\in A$, 设$g:s(b)\to F(b)$是保序双射, 由(1)知$F(a)=g(a)$对$a<b$成立. 因此有$\{x\in \bar{\alpha}\vert x<F(b)\}=\{F(a)\vert a<b\}=F(s(b))=g(s(b))=F(b)$.
\par 由(i)(ii)引理得证. 
\par 回到原题, 设$a\in X$, 若$\forall x<a$, 存在序数与$s(x)$相似, 则引理给出存在序数与$s(a)$相似. 由超限归纳法原理, $\forall a\in X$, 存在序数与$s(a)$相似. 由此引理又给出存在序数与$X$相似, 得证. 
\end{proof}

\par 容易验证对良序集$X,Y$, $\operatorname{ord}X=\operatorname{ord}Y\iff X\cong Y$.

\begin{dfn} \textbf{序数的加法}\\
设$\alpha,\beta$是序数, 取不相交的良序集$A,B$, 使得$\alpha=\operatorname{ord}A$, $\beta=\operatorname{ord}B$. 设$A,B$上的序关系分别为$R,S$, 在$A\cup B$上定义序关系$R\cup S \cup (A\times B)$, 则$A\cup B$也为良序集. 定义序数$\alpha,\beta$的和$\alpha+\beta=\operatorname{ord}(A\cup B)$. 
\end{dfn}
\par 若$A\cap B$非空, 令$\Tilde{A}=\{(a,0)\vert a\in A\}$, $\Tilde{B}=\{(b,1)\vert b\in B\}$, 其序关系分别与第一分量在$A,B$中的序关系一致, 则$\Tilde{A}, \Tilde{B}$是不相交的良序集. 如上定义的良序集$A\cup B$称为良序集$A,B$的有序和(ordinal sum). 可以验证序数的和不依赖良序集$A,B$的选取. 

\begin{pro}\textbf{序数加法的性质}\\
设$\alpha,\beta, \gamma$是序数, 则有
\begin{enumerate}
\item $\alpha+0=\alpha$, $0+\alpha=\alpha$;
\item $\alpha+1=\alpha^+$;
\item 结合律: $(\alpha+\beta)+\gamma = \alpha+(\beta+\gamma)$.
\item 左消去律: $\alpha+\beta=\alpha+\gamma \Rightarrow \beta=\gamma$.
\end{enumerate}
\end{pro}
\begin{proof}
只证左消去律. 取不相交的良序集$A,B,C$, 满足$\operatorname{ord}A=\alpha$, $\operatorname{ord}B=\beta$, $\operatorname{ord}C=\gamma$, 则存在有序和间的保序双射$f:A\cup B \to A\cup C$. 由超限归纳法可得$\forall a\in A: f(a)=a$, 故$f(A)=A$. 因此$f|_B$是$B$到$C$的保序双射, 这表明$B\cong C$, 故$\beta=\gamma$. 
\end{proof}
序数加法不满足交换律, 例如$w+1=w^+$, 但$1+w=w$. 

\begin{dfn} \textbf{序数和的推广}\\
设$\{\alpha_i\}_{i\in I}$是一族序数, 且指标集$I$是良序集. 取两两不相交的一族良序集$\{A_i\}_{i\in I}$, 使得$\alpha_i=\operatorname{ord}A_i$. 在$\bigcup_{i\in I} A_i$上定义序关系如下: 对$a\in A_i$, $b\in A_j$, 若$i\neq j$, $a<b$当且仅当在$I$中有$i<j$; 若$i=j$, $a<b$当且仅当在$A_i$中有$a<b$. 此时$\bigcup_{i\in I} A_i$也为良序集. 定义$\sum_{i\in I} \alpha_i=\operatorname{ord}\left(\bigcup_{i\in I} A_i\right)$. 
\end{dfn}
\par 若$\{A_i\}_{i\in I}$不满足两两不相交的条件, 令$\Tilde{A}_i=\{(a,i)\vert a\in A_i\}$, 其序关系与第一分量在$A_i$中的序关系一致, 则$\{\Tilde{A}_i\}_{i\in I}$两两不相交. 可以验证上述定义不依赖$A_i$的选取. 

\begin{dfn} \textbf{序数的乘法}\\
设$\alpha,\beta$是序数, 任取良序集$A,B$, 使得$\alpha=\operatorname{ord}A$, $\beta=\operatorname{ord}B$. 在$A\times B$上定义序关系如下: $(a,b)<(c,d)$当且仅当$b<d$或$b=d\land a<c$, 则$A\times B$是良序集. 定义序数$\alpha,\beta$的积$\alpha\beta=\operatorname{ord}(A\times B)$. 
\end{dfn}
如上定义的良序集$A\times B$称为良序集$A,B$的有序积(ordinal product). 可以验证序数的积不依赖良序集$A,B$的选取. 

\begin{pro}\textbf{序数乘法的性质}\\
设$\alpha,\beta, \gamma$是序数, 则有
\begin{enumerate}
\item $\alpha 0=0$, $0\alpha =0$;
\item $\alpha 1=\alpha$, $1\alpha =\alpha$;
\item 结合律: $(\alpha\beta)\gamma=\alpha(\beta\gamma)$;
\item 左分配律: $\alpha(\beta+\gamma)=\alpha\beta+\alpha\gamma$.
\end{enumerate}
\end{pro}
序数乘法不满足交换律、右分配律, 例如$2w=(1+1)w=w$, 但$w2=w+w\neq w$. 

\section*{阅读材料}
\par \cite{NST}包含本章内容的概论.

\section*{问题}
\newcounter{probpar} 

\stepcounter{probpar}
\par \textbf{\theprobpar .} 设$R, S$是集合$X$上的关系. (1) 若$R$是偏序关系, $\forall x,y\in X: xSy \iff (xRy \land x\neq y)$. 证明$S$满足反自反性、禁对称性、传递性. (2) 若$S$满足反自反性、反对称性、传递性, $\forall x,y\in X: xRy \iff (xSy \lor x=y)$, 证明$R$是偏序关系. 

\stepcounter{probpar}
\par \textbf{\theprobpar .} \textbf{(偏序集的直积)} 设$(X_i,\le_i),1\le i\le n$是$n$个偏序集, 在$X=X_1\times X_2 \times \cdots\times X_n$上定义关系$T:(x_1,x_2,\cdots,x_n)T(y_1,y_2,\cdots,y_n)$ 当且仅当$x_i\le_i y_i$ 对$\forall 1\le i\le n$成立. 证明: $(X_1\times X_2 \times \cdots\times X_n, T)$ 是偏序集, 称为这$n$个偏序集的直积.

\stepcounter{probpar}
\par \textbf{\theprobpar .} 设$X$是非空有限偏序集, 证明: (1) $X$有极小元、极大元. (2)若$X$是全序集, $X$有最小元、最大元. 

\stepcounter{probpar}
\par \textbf{\theprobpar .} 设$X$是非空有限集, $|X|=n$, $\le$是$X$上的全序关系. 证明: (1)存在$X$的全排列$(x_1,\dots,x_n)$, 满足$x_i\le x_j \iff i\le j$. (2) $X$上的全序关系有$n!$个. 

\stepcounter{probpar}
\par \textbf{\theprobpar .} 设偏序集$X$的偏序关系为$R$, 覆盖关系为$S$, 记$T=\bigcup_{n=1}^\infty S^n$, 其中$S^n$为$n$个$S$的复合, $I=\{(x,x)\vert x\in X\}$. 证明: (1) $I\cup T\subseteq R$. (2)若$X$是有限集, $I\cup T=R$.

\stepcounter{probpar}
\par \textbf{\theprobpar .} \cite{NST} 证明以下命题与Zorn引理等价:
\par (1) 设$(X, \le)$是偏序集, 考虑$\mathcal{X}=\{C\subseteq X\vert C\text{是链}\}$在子集关系下构成的偏序集, $\mathcal{X}$有极大元. 
\par (2) 设$(X, \le)$是偏序集, 考虑$\mathcal{X}=\{C\subseteq X\vert C\text{是链}\}$在子集关系下构成的偏序集, 则$\forall C\in \mathcal{X}$, 存在$\mathcal{X}$中的极大元$C_m$, $C\subseteq C_m$. 
\par (3) 设$(X,\le)$是偏序集, 若对任意链$A\subseteq X$, $X$中存在$A$的上确界, 则$X$有极大元. 

\stepcounter{probpar}
\par \textbf{\theprobpar .} 用Zorn引理证明选择公理. 

\stepcounter{probpar}
\par \textbf{\theprobpar .} 用良序定理证明选择公理.


\stepcounter{probpar}
\par \textbf{\theprobpar .} \cite{NST}设$X$是全序集, 证明存在良序集$S\subseteq X$, 满足$\forall x\in X\exists a\in S: x\le a$. 

\stepcounter{probpar}
\par \textbf{\theprobpar .} \cite{NST}设$X$是全序集, 证明$X$是良序集当且仅当$\forall a\in X$, $s(a)=\{x\in X\vert x<a\}$是良序集. 

\stepcounter{probpar}
\par \textbf{\theprobpar .} \cite{NST} 设$X$是全序集, 且$X$的任一可数子集是良序集, 证明$X$是良序集. 

\stepcounter{probpar}
\par \textbf{\theprobpar .} \cite{IC} 设$R$是有限集$X$上的偏序关系. 构造$X$上的全序关系$S$, 满足$R\subseteq S$. 

\stepcounter{probpar}
\par \textbf{\theprobpar .} \textbf{(Szpilrajn扩张定理)} 设$R$是$X$上的偏序关系, 证明存在$X$上的全序关系$S$, 满足$R\subseteq S$. 

\stepcounter{probpar}
\par \textbf{\theprobpar .} \cite{NST}设$X,Y$是良序集, 存在双射$f:X\to Y$, 满足$\forall a,b\in X$, 在$X$中$a\le b$当且仅当在$Y$中$f(a)\le f(b)$. 证明这样的双射唯一. 

\stepcounter{probpar}
\par \textbf{\theprobpar .} 设$\alpha$是非零序数, 证明$\alpha$是极限序数当且仅当对任意序数$\beta<\alpha$, 存在序数$\gamma$, 满足$\beta<\gamma<\alpha$. 

\stepcounter{probpar}
\par \textbf{\theprobpar .} \cite{NST}设$\alpha,\beta$是序数, 证明$\alpha<\beta$当且仅当存在非零序数$\gamma$, $\alpha+\gamma=\beta$. 

\section*{解答}
\newcounter{solpar} 

\stepcounter{solpar}
\par \textbf{\thesolpar .} (1) 反自反性显然. 禁对称性: 若$xSy\land ySx$, $xRy \land yRx \land x\neq y$, 与$R$的反对称性矛盾. 传递性: 若$xSy\land ySz$, 则$xRy \land yRz$, $xRz$. 若$x=z$, 由$zRy \land yRz$推出$y=z$, 与$ySz$矛盾, 因此$x\neq z$, 故$xSz$. (2) 自反性显然. 反对称性: 若$xRy \land yRx$, 反设$x\neq y$, 则$xSy \land ySx$, 推出$x=y$, 矛盾. 传递性: 若$xRy \land yRz$, 分$x=y=z$, $x=y\neq z$, $x\neq y=z$, $x\neq y\neq z$四种情况讨论得$xRz$. 

\stepcounter{solpar}
\par \textbf{\thesolpar .} 证明略. 

\stepcounter{solpar}
\par \textbf{\thesolpar .} (1)以极大元为例. 对$n$归纳. 当$n=1$, $X$唯一的元素是极大元. 设命题对元素个数为$n$的情形成立, 考虑$n+1$的情形. 取$a\in X$, 则$|X-\{a\}|=n$, 从而$X-\{a\}$有极大元$m$. 若$m\le a$, 由传递性$a$为$X$的极大元; 否则$m$也为$X$的最大元. (2)全序集的极大元是最大元.  

\stepcounter{solpar}
\par \textbf{\thesolpar .} (1)取$x_1$为$X$的最小元, $x_2$为$X-\{x_1\}$的最小元, ..., $x_n$为$X-\{x_1,\dots, x_{n-1}\}$的最小元. (2)基于(1)定义全序集合到全排列集合的映射, 则$x_i\le x_j \iff i\le j$给出逆映射, 故全序个数等于全排列数$n!$. 

\stepcounter{solpar}
\par \textbf{\thesolpar .} (1) 显然$I\subseteq R$, $S\subseteq R$. 对$n$归纳, 设$S^n\subseteq R$, 对$(x,y)\in S^{n+1}$, 存在$z\in X$, $(x,z)\in S^n$, $(z,y)\in S$. 由传递性, $(x,y)\in R$, 故$S^{n+1}\subseteq R$. (2) 对$X$的元素个数归纳, $|X|=0,1$的情形显然. 假定当$|X|\le k$时有$R\subseteq I\cup T$. 考虑$|X|=k+1$的情形, 任取$x\le y$, 不妨设$x\neq y$. 令$A=\{z\in X\vert x<z<y\}$. 若$A=\varnothing$, 则$(x,y)\in S$. 否则, $A\cup \{y\}$在$X$的偏序关系下也构成有限偏序集, 且$|A\cup \{y\}|\le k$. 取$A$的一个极小元$z$, 则$(x,z)\in S$, 由归纳假设$(z,y)\in T$, 故$(x,y)\in T$. 命题得证. 
\par \emph{注:} (2)对无限集不成立, 如有理数集$\mathbb{Q}$的覆盖关系是空集. 

\stepcounter{solpar}
\par \textbf{\thesolpar .} 记Zorn引理为(Z), 其证明过程表明(1)$\to$(Z); 由$\mathcal{X}$非空知(2)$\to$(1), 又(Z)$\to$(3)显然. 只需证(3)$\to$(2). 对任意给定的$C\in \mathcal{X}$, 考虑$\mathcal{X}_C=\{A\in \mathcal{X}\vert C\subseteq A\}$在子集关系下构成的偏序集, 对任意链$\mathcal{C} \subseteq \mathcal{X}_C$, $\bigcup_{A\in \mathcal{C}} A$是链. 事实上, 对$x_1,x_2\in \bigcup_{A\in \mathcal{C}} A$, 存在$A_1, A_2\in \mathcal{C}$, $x_1\in A_1$, $x_2\in A_2$. 由$\mathcal{C}$是链, $A_1, A_2$有包含关系, 不妨设$x_2\in A_1$, 由$A_1$是链, 知$x_1,x_2$可比. 进一步有$\bigcup_{A\in \mathcal{C}} A\in \mathcal{X}_C$, 且是$\mathcal{C}$的上确界. 因此, $\mathcal{X}_C$有极大元$C_m$; 对$C'\in \mathcal{X}$, 若$C_m\subseteq C'$, 则由$C\subseteq C'$知$C'\in \mathcal{X}_C$, 从而$C'=C_m$. 故$C_m$也是$\mathcal{X}$的极大元. 

\stepcounter{solpar}
\par \textbf{\thesolpar .} \cite{NST} 我们证明(C)的等价命题: 设$X$是非空集合, 存在函数$f:\mathcal{P}(X)-\{\varnothing\}\to X$, 满足$\forall A\in \mathcal{P}(X)-\{\varnothing\}$, $f(A)\in A$. 考虑集合$\mathcal{F}=\{f\in \mathcal{P}(\mathcal{P}(X)\times X)\vert f\text{是函数且}\forall A\in \text{dom}f: f(A)\in A\}$. 显然$\text{dom}f\subseteq \mathcal{P}(X)-\{\varnothing\}$. 用$f\le g$表示$g$是$f$的延拓, 则$(\mathcal{F},\le)$是偏序集. 设$\mathcal{C}\subseteq \mathcal{F}$是链, 则$\bigcup_{f\in \mathcal{C}} f\in \mathcal{F}$是$\mathcal{C}$的上界. 由Zorn引理, $\mathcal{F}$存在极大元$f_m$. 断言$\text{dom}f_m=\mathcal{P}(X)-\{\varnothing\}$. 若不然, 取$X$的非空子集$A\notin \text{dom}f_m$, 并取$a\in A$, 则$f_m \cup \{(A,a)\}\in \mathcal{F}$且$f_m < f_m \cup \{(A,a)\}$, 与$f_m$是极大元矛盾. 

\stepcounter{solpar}
\par \textbf{\thesolpar .} 我们证明(C)的等价命题: 设$X$是非空集合, 存在函数$f:\mathcal{P}(X)-\{\varnothing\}\to X$, 满足$\forall A\in \mathcal{P}(X)-\{\varnothing\}$, $f(A)\in A$. 由良序定理, 可设$X$是良序集, 此时令$f=\{(A,a)\in \mathcal{P}(X)\times X\vert a\text{为}A\text{的最小元}\}$, 即满足要求. 
\par \emph{注}: 本题表明选择公理、Zorn引理、良序定理三者等价. 


\stepcounter{solpar}
\par \textbf{\thesolpar .} 类似良序定理的证明, 令$\mathcal{W}=\{A\subseteq X\vert A\text{是良序集}\}$, $\mathcal{W}$在延续关系下构成偏序集. 设$\mathcal{C}\subseteq \mathcal{W}$是链, 则存在$U=\bigcup_{C\in \mathcal{C}} C$是$C$的上界. 由Zorn引理, $\mathcal{W}$存在极大元$S$. 若存在$x\in X$, $\forall a\in S: x>a$, 则$S\cup \{a\}$是良序集且是$S$的延续, 与$S$为极大元矛盾. 

\stepcounter{solpar}
\par \textbf{\thesolpar .} 必要性显然. 充分性: 设非空集合$A\subseteq X$, 则存在$a\in A$. 若$A\cap S(a)$非空, 则存在$A\cap S(a)$的最小元$x$, 对$y\in A-S(a)$, 有$x<a\le y$, 因此$x$也是$A$的最小元; 若$A\cap S(a)= \varnothing$, 则$a$是$A$的最小元. 

\stepcounter{solpar}
\par \textbf{\thesolpar .} 反设$X$不是良序集, 则存在非空集合$A\subseteq X$, $A$没有最小元. 若$A$是可数集, 已构成矛盾. 若不然, 由选择公理, 存在函数$f:\mathcal{P}(A)-\{\varnothing\}\to A$, 满足对任意非空集合$B\subseteq A, f(B)\in B$. 递归定义函数$u:\mathbb{N}\to A$如下: $u(0)=f(A)$; $\forall n\in \mathbb{N}$, 记$P_n=\{x\in A\vert x<u(n)\}$, 由$A$是全序集且不存在最小元, $P_n$非空, 令$u(n+1)=f(P_n)$. 显然$u(\mathbb{N})$是可数集, 但$\forall n\in \mathbb{N}: u(n+1)<u(n)$, 故其没有最小元, 与$u(\mathbb{N})\subseteq A$是良序集矛盾. 

\stepcounter{solpar}
\par \textbf{\thesolpar .} \cite{IC} 记$|X|=n$. 取$X$的一个极小元, 记为$x_1$; 取$X-\{x_1\}$的一个极小元, 记为$x_2$; ... 取$X-\{x_1,\dots,x_{n-1}\}$的一个极小元, 记为$x_n$. 定义$X$上的全序$S$如下: $x_i S  x_j \iff i\le j$. 为验证$R\subseteq S$, 反设存在$x_j<x_i, i<j$, 则与$x_i$是$X-\{x_1,\dots,x_{i-1}\}$的极小元矛盾. 

\stepcounter{solpar}
\par \textbf{\thesolpar .} 设$\mathcal{R}=\{S\in \mathcal{P}(X\times X)\vert R\subseteq S \land S\text{是偏序关系}\}$, 则$\mathcal{R}$在子集关系下构成偏序集. 设$\mathcal{C}\subseteq \mathcal{R}$是链, 则$\bigcup_{S\in \mathcal{C}}S\in \mathcal{R}$是$\mathcal{C}$的上界. 由Zorn引理, $\mathcal{R}$中存在极大元$S_M$, 我们证明$S_M$是全序关系. 若不然, 设$(x,y)\notin S_M \land (y,x)\notin S_M$, 显然有$x\neq y$. 令$W=\{(z,w)\vert (z,x)\in S_M \land (y,w)\in S_M\}$, $S_M'=S_M\cup W$, 由$(x,y)\in W$知$S_M\subset S_M'$. $S_M'$显然满足自反性. 设$(z,w)\in W$, 若$(w,z)\in S_M$或$(w,z)\in W$, 由$S_M$的传递性均有$(y,x)\in S_M$, 矛盾. 故$S_M'$仍满足反对称性. 设$(u,v)\in S_M'$, $(v,w)\in S_M'$, 分三种情况：(1)$(u,v),(v,w)\in S_M$, 显然有$(u,w)\in S_M$; (2)$(u,v),(v,w)$恰有一个在$W$中, 则$(u,x)\in S_M$, $(y,w)\in S_M$, 得$(u,w)\in W$; (3)$(u,v),(v,w)\in W$, 此时有$(y,v)\in S_M$, $(v,x)\in S_M$, 这推出$(y,x)\in S_M$, 矛盾. 故$S_M'$仍满足传递性. 综上所述, $S_M'$是偏序集, 这与$S_M$是极大元矛盾. 

\stepcounter{solpar}
\par \textbf{\thesolpar .} \cite{NST} 设$g$也为保持序关系的双射, 则$h=f^{-1}\circ g$是$X$到自身且保持序关系的双射. 由引理\ref{lem:wossim}, $\forall a\in X: a\le f^{-1}(g(a))$, 亦即$f(a)\le g(a)$. 而$g^{-1}\circ f$也是$X$到自身且保持序关系的双射, 故同理可得 $\forall a\in X: g(a)\le f(a)$. 因此$g=f$. 

\stepcounter{solpar}
\par \textbf{\thesolpar .} 必要性: 设$\alpha$是极限序数, 对$\beta<\alpha$, 有$\beta^+\neq \alpha$. 若$\alpha<\beta^+$, 则$\alpha\in \beta^+$, $\alpha \in \beta$或$\alpha=\beta$, 均与$\beta<\alpha$矛盾. 由序数的三歧性, 必有$\beta^+<\alpha$. 充分性: 反证法, 若$\alpha=\beta^+$, 则$\forall \gamma>\beta$, 有$\beta \in \gamma$, $\beta \subseteq \gamma$. 这表明$\alpha \subseteq \gamma$, 与存在$\beta<\gamma<\alpha$矛盾. 

\stepcounter{solpar}
\par \textbf{\thesolpar .} 必要性: 由$\alpha<\beta$, $\beta$是$\alpha$的延续. 记$A=\beta-\alpha$, 则$A$是非空良序集, 且根据序数加法的定义有$\alpha+\operatorname{ord}A=\beta$, 则$\gamma=\operatorname{ord}A$满足条件. 充分性: 反证法, 若$\beta\le \alpha$, 由必要性的证明知存在序数$\delta$, $\alpha=\beta+\delta$. 因此有$\beta+\delta+\gamma=\beta$. 由左消去律, $\delta+\gamma=0$, 与$\gamma\neq 0$矛盾. 